\section{Climate: Exposure Without Influence}
\label{sec:climate}

Globalisation did not cause cyclones or sea-level rise, and this paper does
not claim it did. The connection is an \emph{asymmetry}: the places most
exposed to climate hazards contribute least to emissions and have the least
say over global responses. Remoteness and smallness sharpen that asymmetry.
A single cyclone can damage a large share of a small island's housing stock and
infrastructure at once, and recovery depends on imported materials, foreign
contractors and international finance---all subject to the distance tariff.

\begin{figure}[t]
\centering
\begin{tikzpicture}
\begin{axis}[edge, height=0.6\linewidth,
    xmode=log, xmin=0.03, xmax=60,
    xtick={0.1,0.3,1,3,10,30}, xticklabels={0.1,0.3,1,3,10,30},
    xlabel={CO$_2$ emissions per person, tonnes (log scale)},
    ymode=log, ymin=0.005, ymax=100,
    ytick={0.01,0.1,1,10,100}, yticklabels={0.01,0.1,1,10,100},
    ylabel={Disaster displacements per 1,000 people, yearly (log)},
    legend pos=south east]
  \addplot[dotsBase, only if={sids}{0}] table[x=co2pc,y=displaced_per_k] {\datadir/plot_countries.csv};
  \addlegendentry{Other economies}
  \addplot[dotsSIDS, only if={sids}{1}] table[x=co2pc,y=displaced_per_k] {\datadir/plot_countries.csv};
  \addlegendentry{Small island states}
  \addplot[labelled=lab_climate] table[x=co2pc,y=displaced_per_k] {\datadir/plot_countries.csv};
\end{axis}
\end{tikzpicture}
\caption{\takeaway{Exposure without influence.} Small island states account
  for \sidsCOshare\% of the world's CO$_2$ emissions (and \sidsPopShare\% of its
  people) yet the median island state sees \dispSIDSRatio{} times as many
  disaster displacements per head as the median other country. \nDispTopTenSIDS{} of the ten most-displaced
  economies are islands. Emissions \windowStart--\windowEnd{} average; displacements
  2008--2023 average (years without a report counted as zero; countries without any report omitted).
  \source{World Bank WDI (EDGAR emissions; IDMC displacement data).}}
\label{fig:climate}
\end{figure}


\Cref{fig:climate} sets per-person carbon emissions against a measure of
realised exposure: new disaster displacements per 1,000 residents, averaged
over 2008--2023 (IDMC data via the World Bank). Small island states account for \sidsCOshare\% of
global CO$_2$ emissions and \sidsPopShare\% of world population; the Pacific
island states for just \pacificCOshare\%. Median emissions are
\coPacificMedian{} tonnes a head in the Pacific and \coSIDSMedian{} in all
island states, against \coHICMedian{} in high-income economies. Yet the
median island state records \dispSIDSRatio{} times as many disaster
displacements per head as the median other economy, and \nDispTopTenSIDS{}
of the ten most-affected countries are islands; \dispTopName{} tops the list
with \dispTopValue{} displacements per 1,000 people a year. Low-lying atolls
face the slow-onset threat too: \lowElevMarshall\% of Marshall Islanders and
\lowElevTuvalu\% of Tuvaluans live less than five metres above sea level.

Exposure alone is not destiny either: \lowElevNetherlands\% of the Dutch also
live below five metres, and the Netherlands is rich and safe because it can
afford to engineer its coast. The asymmetry that matters is between exposure
and the \emph{capacity} to adapt---money, engineers, insurance and
negotiating weight---and capacity is exactly what the recurring tax erodes.
